Here we provide further detail on the model described in Eq.~\eqref{eq:neg_binom_mixed_effects}.

For some elements of the parameters $\phi$, we have a prior that informs point estimation in MAP fashion. 
The random effects are assumed to follow a multivariate normal prior
\begin{equation}
    \begin{pmatrix} b_{0s} \\ b_{1s} \end{pmatrix} \sim \mathcal{N}\left(\begin{pmatrix} 0 \\ 0 \end{pmatrix}, \mathbf{\Sigma}\right)
\end{equation}
where the covariance matrix $\mathbf{\Sigma}$ is parameterized as:
\begin{equation}
    \mathbf{\Sigma} = \begin{pmatrix} 
    \sigma_0^2 & \rho\sigma_0\sigma_1 \\ 
    \rho\sigma_0\sigma_1 & \sigma_1^2 
    \end{pmatrix}
\end{equation}
Here, $\sigma_0$ is the standard deviation of the random intercepts, $\sigma_1$ is the standard deviation of the random slopes, and $\rho$ is the correlation between the random intercepts and slopes, where $-1 \leq \rho \leq 1$. We pack these hyperparameters into a separate vector $\eta = \{\sigma_0, \sigma_1, \rho \}$.

\textbf{Transforming to unconstrained parameters.}
To ensure parameter constraints are satisfied throughout gradient descent optimization, we employ invertible transformations that map constrained domains to unconstrained spaces. We reparameterize the correlation coefficient $\rho \in (-1, 1)$ using $u \gets \text{arctanh}(\rho)$, allowing $u$ to be optimized over $\mathbb{R}$ while ensuring $\rho = \tanh(u)$ remains within its valid bounds. For the strictly positive parameters $\sigma_0, \sigma_1 > 0$, we optimize unconstrained parameters $\xi_0, \xi_1 \in \mathbb{R}$ and apply the softplus transformation $\sigma_i \gets \ln(1 + e^{\xi_i})$, which guarantees positivity while maintaining smoothness. Similarly, for the probability parameter $q \in (0, 1)$, we optimize an unconstrained parameter $\zeta \in \mathbb{R}$ and apply the sigmoid transformation $q \gets 1/(1 + e^{-\zeta})$. During gradient descent, we optimize these unconstrained parameters, and transformed to the constrained version during forward sampling or pdf evaluation of the model.

\textbf{MAP estimation.}
Together, the parameters $\phi$ and hyperparameters $\eta$ comprise this model. 
Both are point estimated to maximize the MAP objective. 
Thus, what is marked in the paper as NLL optimization is really best viewed as MAP or penalized NLL estimation, where the loss is
\begin{equation}
    \mathcal{J}(\phi, \eta) = - \log p(\bm{y}_{1:T} |\bm{\phi}) - \log p( \phi | \eta)
\end{equation}
Similarly, when we fit this model with DAML, the above loss is what is minimized subject to the BPR constraint.